%% Forest Ecology and Management — Elsevier submission
%% Document class: elsarticle (review mode = double spacing + line numbers)
%% Reference style: numbered [elsarticle-num-names]

\documentclass[review,12pt,a4paper]{elsarticle}

\usepackage{lineno}          % Line numbers (required for review)
\usepackage{natbib}          % Bibliography
\usepackage{amsmath}         % Math environments
\usepackage{amssymb}
\usepackage{booktabs}        % Professional tables (\toprule, \midrule, \bottomrule)
\usepackage{longtable}       % Tables spanning pages
\usepackage{tabularx}        % Auto-width columns
\usepackage{array}           % Advanced column types
\usepackage{multirow}        % Row spanning in tables
\usepackage{graphicx}        % Figures
\usepackage{xcolor}          % Colour support
\usepackage{url}             % URL formatting
\usepackage{hyperref}        % Hyperlinks (load last)
\usepackage{cleveref}        % Smart cross-references

\hypersetup{
  colorlinks = true,
  linkcolor  = blue!70!black,
  citecolor  = blue!70!black,
  urlcolor   = blue!70!black
}

\modulolinenumbers[5]
\linenumbers

%% ============================================================
\begin{document}

\begin{frontmatter}

%% ── TITLE ───────────────────────────────────────────────────
\title{Forest Stratification and Vertical Zonation across Environmental
Gradients in Bhutan: Diversity Patterns, Community Assembly, and
Long-Term Greenness Trends from the National Forest Inventory}

%% ── AUTHORS AND AFFILIATIONS ────────────────────────────────
\author[aff1]{Wangdi\corref{cor1}}
\ead{wangdiues@gmail.com}

\author[aff2]{Sangay Chedup}
\ead{sangaychedup16@gmail.com}

\cortext[cor1]{Corresponding author. Tel.: +975-17506599.
  ORCID: \href{https://orcid.org/0009-0007-7726-1742}{0009-0007-7726-1742}}

\affiliation[aff1]{organization={Forest Resources Planning \& Management Division,
  Department of Forests and Park Services,
  Ministry of Energy and Natural Resources,
  Royal Government of Bhutan},
  city={Thimphu},
  country={Bhutan}}

\affiliation[aff2]{organization={Divisional Forest Office, Mongar,
  Department of Forests and Park Services,
  Ministry of Energy and Natural Resources,
  Royal Government of Bhutan},
  city={Thimphu},
  country={Bhutan}}

%% ── ABSTRACT ────────────────────────────────────────────────
\begin{abstract}
Bhutan's forests span one of the steepest unbroken altitudinal gradients
in the Eastern Himalayas (113--5{,}469~m), yet a country-wide, multi-method
characterisation of their vertical stratification remains absent.
Using 1{,}942 systematically distributed National Forest Inventory (NFI)
plots recording 2{,}195 taxa across 12 forest types, we simultaneously
quantified alpha diversity, beta-diversity turnover, environmental
drivers of community assembly, indicator species, co-occurrence network
structure, structural complexity, and long-term vegetation greenness
trends. Species richness declined monotonically from lowland forest
(24.8~$\pm$~9.0 species per plot at 500--1{,}000~m) to alpine scrub
(8.9~$\pm$~5.6 above 4{,}000~m; Kruskal-Wallis $H = 630.8$,
$\eta^2 = 0.322$), with substantial variation among forest types.
PERMANOVA confirmed significant compositional differentiation among
forest types (pseudo-$F = 17.03$, $R^2 = 0.089$, $p = 0.001$;
999~permutations), and ecological canonical correspondence analysis
(CCA) identified cold-season minimum temperature (\textit{bio6}) and
elevation as the primary community-assembly drivers
(CCA1: 2.45\% explained; full-model adj-$R^2 = 4.22$\%;
pure climate fraction = 2.90\%).
Indicator species analysis recovered 416 significant taxon--habitat
associations, with the highest habitat fidelity in conifers
(\textit{Pinus roxburghii}, IndVal = 0.794;
\textit{Picea spinulosa}, 0.463) and \textit{Rhododendron} species in the
alpine zone (\textit{R.\ anthopogon}, 0.283).
Co-occurrence network analysis recovered 1{,}083 species nodes and
26{,}875 edges structured into five community modules
(modularity $Q = 0.287$; null-model SES = 0.615, $p = 0.242$),
with the three dominant modules corresponding to lowland broadleaf,
montane mixed, and subalpine--alpine assemblages.
The Stratification Complexity Index was highest in subtropical and warm
broadleaved forests, lowest in alpine scrub, and robust to component
weighting (all leave-one-out Spearman $r \geq 0.983$).
Pixel-level MODIS EVI trend analysis (2000--2024, {\raise.17ex\hbox{$\scriptstyle\sim$}}250~m resolution,
Mann-Kendall) revealed that 29.9\% of Bhutan's forested area experienced
significant long-term greening after Benjamini--Hochberg false discovery
rate correction (FDR~$\leq 0.05$; 210{,}180~pixels; 11{,}655~km$^2$),
with only 0.7\% showing significant browning.
Greening intensity was strongest at low elevations
($<$500~m: 70.6\% of plots significantly greening) and weakest in the
mid-montane zone (3{,}500--4{,}000~m: 18.9\%).
EVI trend slope correlated positively with plot-level species richness
(Pearson $r = 0.107$, $R^2 = 0.012$, $p < 0.001$), though explained
variance was low and causal inference is not supported.
Together, these findings establish a reproducible, open-source baseline
for long-term biodiversity monitoring and climate-adaptive forest
management in the Eastern Himalayas.
\end{abstract}

%% ── KEYWORDS ────────────────────────────────────────────────
\begin{keyword}
altitudinal zonation \sep
vertical stratification \sep
species richness \sep
canonical correspondence analysis \sep
indicator species \sep
co-occurrence network \sep
stratification complexity index \sep
MODIS EVI \sep
Eastern Himalayas \sep
National Forest Inventory
\end{keyword}

\end{frontmatter}

%% ============================================================
\section{Introduction}
\label{sec:intro}

Mountain forests occupy approximately 28\% of the world's forested land
area yet shelter a disproportionate share of global plant diversity, owing
to the compression of multiple climatic and edaphic regimes across short
horizontal distances \citep{korner2012}. Altitudinal gradients function as
natural experiments in community ecology: as elevation increases, temperature
declines, growing seasons shorten, soil development is truncated, and the
physiological cost of cold tolerance progressively filters the species pool,
generating predictable turnover in both composition and forest structural
complexity \citep{rahbek1995}. The Hindu Kush--Himalayan arc is among the
world's most significant biodiversity hotspots \citep{myers2000}, harbouring
an estimated 25{,}000 plant species and containing several of the world's
steepest biotic gradients --- yet the majority of its floristic diversity
remains incompletely characterised at the landscape scale. Superimposed on
this baseline, ongoing climate change is altering altitudinal vegetation zones
in mountain ecosystems faster than in lowland counterparts \citep{pepin2015},
advancing treelines, shifting species range boundaries, and modifying forest
structural complexity in ways that are only beginning to be captured by
systematic inventory data. Establishing rigorous, reproducible baselines of
forest composition, vertical stratification, and long-term vegetation dynamics
is therefore both scientifically urgent and practically essential for
evidence-based conservation management in Himalayan nations.

Bhutan occupies a pivotal position within this context. Spanning
26.7$^\circ$--28.4$^\circ$N and 88.7$^\circ$--92.2$^\circ$E in the
Eastern Himalayas, the country encompasses an unbroken elevational transect
from 113~m in the subtropical southern foothills to 7{,}570~m at glaciated
peaks --- a vertical compression of five climatic zones (subtropical, warm
temperate, cool temperate, subalpine, and alpine) within fewer than 150~km
of horizontal distance. Under a constitutional mandate requiring at least
60\% national forest cover, Bhutan retains approximately 71\% of its land
area as forest, distributed across twelve officially recognised forest types
from Chirpine (\textit{Pinus roxburghii}) and subtropical broadleaved
communities in the south to Fir (\textit{Abies densa}), Spruce
(\textit{Picea spinulosa}), Blue Pine (\textit{Pinus wallichiana}), and
alpine \textit{Rhododendron}--\textit{Juniperus} scrub in the north. This
configuration --- extreme altitudinal relief, high forest cover, low direct
anthropogenic pressure, and a formally legislated conservation mandate ---
makes Bhutan an exceptional natural laboratory for investigating
altitude-driven forest stratification and the multi-scale environmental
filters governing plant community assembly. Critically, Bhutan's first
systematic National Forest Inventory (NFI), comprising 1{,}942 permanent
sample plots distributed across the national forest estate, provides for the
first time a country-wide, spatially representative dataset with which to
characterise these patterns empirically.

Despite its ecological significance, Bhutan's forest vertical zonation has
not previously been characterised at the national scale using multi-method
quantitative analysis. Three specific knowledge gaps motivate the present
study. First, while altitudinal species richness gradients are well documented
in parts of the Himalayas \citep{grytnes2002}, the simultaneous quantification
of alpha diversity, layer-specific richness, and community assembly drivers
using field inventory data across the full elevational range of an Eastern
Himalayan country has not been attempted. Second, the topological structure of
plant species co-occurrence networks in Himalayan forests --- including the
degree of community modularity relative to null expectations --- remains
largely undescribed, with most co-occurrence network analyses focusing on
lowland tropical or temperate systems. Third, spatially-explicit, long-term
satellite-derived vegetation greenness trends have not been integrated with
NFI-derived diversity and structural complexity data at the national scale in
Bhutan, leaving a critical evidence gap for REDD+ monitoring and climate
adaptation planning.

This study addresses these gaps through a comprehensive, integrated analysis
of Bhutan's NFI dataset, with six specific objectives:
\begin{enumerate}
  \item to quantify alpha diversity (species richness, Shannon entropy, and
        layer-specific richness for trees, shrubs, and herbs) across the full
        elevational gradient and among the 12 forest types;
  \item to characterise community beta-diversity turnover using non-metric
        multidimensional scaling (NMDS) and permutational multivariate analysis
        of variance (PERMANOVA), and to identify primary environmental drivers
        of community assembly by ecological CCA with variance partitioning;
  \item to identify indicator taxa diagnostic of each forest type using the
        IndVal method \citep{dufresne1997};
  \item to describe the topology and modular community structure of the
        national species co-occurrence network relative to a swap-permutation
        null model;
  \item to compute and map a Stratification Complexity Index (SCI)
        integrating multi-layer structural information into a single composite
        metric; and
  \item to assess spatially-explicit 24-year vegetation greenness trends using
        MODIS EVI data (2000--2024, {\raise.17ex\hbox{$\scriptstyle\sim$}}250~m resolution) at both the pixel and
        NFI-plot level, and to link these trends to elevation zone, forest
        type, and species diversity.
\end{enumerate}

The study represents the first integrated, country-scale, multi-method
synthesis of Bhutan's forest vertical zonation and provides three novel
contributions. Analytically, it applies a complete community ecology
workflow --- spanning alpha and beta diversity, constrained ordination,
indicator species analysis, co-occurrence network modelling, structural
complexity indexing, and remote sensing trend analysis --- to a single
national inventory dataset, enabling cross-method synthesis that isolated
analyses cannot achieve. Technically, all analyses are implemented in a
modular, fully reproducible open-source Python pipeline with a fixed random
seed, enabling exact replication and adaptation for future NFI cycles. In
terms of applied relevance, the results directly inform Bhutan's National
Biodiversity Strategy and Action Plan, its REDD+ monitoring and reporting
obligations, and the design of ecologically targeted conservation management
across its twelve forest types.

%% ============================================================
\section{Materials and Methods}
\label{sec:methods}

\subsection{Study area}
\label{sec:study_area}

Bhutan (26.7$^\circ$--28.4$^\circ$N, 88.7$^\circ$--92.2$^\circ$E;
38{,}394~km$^2$) occupies the eastern Himalayas and encompasses one of the
steepest unbroken elevational gradients in the world, spanning from 113~m in
the subtropical foothills to 7{,}570~m at glaciated peaks. Five broadly
recognised climatic zones --- subtropical, warm temperate, cool temperate,
subalpine, and alpine --- are compressed within fewer than 150~km of
horizontal distance. Under a constitutional mandate requiring at least 60\%
forest cover, Bhutan retains approximately 71\% of its land area under forest,
with twelve classified forest types recognised in national land-cover mapping,
ranging from subtropical broadleaved and Chirpine (\textit{Pinus roxburghii})
forests in the south to Fir (\textit{Abies densa}), Spruce
(\textit{Picea spinulosa}), Blue Pine (\textit{Pinus wallichiana}), and alpine
scrub communities in the north.

\subsection{National Forest Inventory data}
\label{sec:nfi_data}

Data were drawn from Bhutan's first systematic National Forest Inventory
(NFI), comprising \textbf{1{,}942 permanent sample plots} distributed on a
systematic grid across forested land. At each plot, field crews recorded the
identity and canopy cover of all vascular plant species in three vertical
layers: Trees ($>$5~m height), Shrubs (0.5--5~m), and Herbs ($<$0.5~m).
Additional plot-level attributes recorded included forest type (one of 12
classes), plot area (ha), soil type, and administrative district (Dzongkhag).
Elevation at each plot centroid ranged from 113 to 5{,}469~m
(mean~$\pm$~SD: 2{,}454~$\pm$~1{,}271~m). Prior to analysis, the species
matrix was cleaned for typographical duplicates, synonyms, and ambiguous
records; plot coordinates were validated against Bhutan's geographic bounding
box. After cleaning, the dataset comprised 2{,}195 unique species taxa across
1{,}942 plots.

\subsection{Environmental variable extraction}
\label{sec:env_extract}

Environmental predictors were extracted at each plot centroid by bilinear
resampling using rasterio (v~$\geq$~1.3). Three predictor groups were
assembled. (i)~\textit{Topography}: elevation (m), slope ($^\circ$), and
aspect ($^\circ$) from the Shuttle Radar Topography Mission (SRTM) 30-m
Digital Elevation Model. (ii)~\textit{Climate}: all 19 CHELSA/WorldClim~v2
bioclimatic variables (\textit{bio1}--\textit{bio19}), representing
temperature and precipitation regimes averaged over 1970--2000. Variables of
particular ecological interest include \textit{bio1} (mean annual temperature,
MAT), \textit{bio6} (minimum temperature of the coldest month),
\textit{bio11} (mean temperature of the coldest quarter), \textit{bio12}
(mean annual precipitation, MAP), and \textit{bio4} (temperature
seasonality). (iii)~\textit{Soil}: FAO soil type class from the Harmonised
World Soil Database. Variables with a Pearson correlation $|r|>0.95$ with
another predictor within the same group were flagged; two computational
artefacts present in an intermediate processing layer were excluded from all
ordination analyses.

\subsection{Alpha diversity}
\label{sec:alpha_div}

Per-plot alpha diversity was quantified using three complementary indices:
species richness ($S$, total count of taxa), Shannon entropy
($H' = -\sum p_i \ln p_i$), and Simpson's diversity
($D = 1 - \sum p_i^2$). Layer-specific richness was computed separately for
Trees, Shrubs, and Herbs. Diversity indices were summarised by nine 500-m
elevation bands ($<$500~m to $>$4{,}000~m) and by forest type. To
statistically justify the 500-m banding scheme, a Kruskal-Wallis $H$ test was
applied to species richness across bands, with eta-squared
($\eta^2 = (H - k + 1)/(n - k)$, where $k$ = number of bands and $n$ =
number of plots) calculated as a rank-based effect size.

\subsection{Beta diversity and community turnover}
\label{sec:beta_div}

Community dissimilarity was quantified using the Bray-Curtis metric on the
log-transformed species abundance matrix ($\log_1(x+1)$). Non-metric
multidimensional scaling (NMDS) was performed in two dimensions using the
SMACOF algorithm (sklearn MDS; \texttt{normalized\_stress=True},
\texttt{max\_iter=300}, \texttt{n\_init=4}, \texttt{random\_state=42}).
With \texttt{normalized\_stress=True}, the \texttt{stress\_} attribute equals
Kruskal formula~1:
\begin{equation}
  \text{stress}_1 = \sqrt{\frac{\sum (d_{ij} - \hat{d}_{ij})^2}{\sum d_{ij}^2}}
  \label{eq:stress}
\end{equation}
where $d_{ij}$ are distances in dissimilarity space and $\hat{d}_{ij}$ are
fitted ordination distances. Stress values below 0.20 are generally
considered acceptable; values above 0.20 indicate the 2D projection does not
faithfully represent the underlying dissimilarity structure. To complement the
NMDS and provide a metric alternative less sensitive to high stress, a
Principal Coordinates Analysis (PCoA; metric MDS) of the same Bray-Curtis
matrix was computed and is provided as Supplementary Fig.~S2.

Compositional differentiation among forest types was tested by PERMANOVA
\citep{anderson2001} using 999 random permutations of plot labels. The
PERMANOVA $R^2$ statistic was computed as:
\begin{equation}
  R^2 = \frac{F \cdot (k-1)}{F \cdot (k-1) + (n-k)}
  \label{eq:permanova_r2}
\end{equation}
where $F$ is the pseudo-$F$ statistic, $k$ is the number of forest-type
groups, and $n$ is the number of plots. Multivariate homogeneity of group
dispersions was assessed by PERMDISP \citep{anderson2006}, implemented as a
Kruskal-Wallis $H$ test across within-group distances to group centroids.

\subsection{Canonical correspondence analysis}
\label{sec:cca}

To identify environmental drivers of community composition, ecological CCA
\citep{terbraak1986} was conducted using \texttt{skbio.stats.ordination.cca}
(scikit-bio~$\geq$~0.6). CCA is a constrained ordination method that relates
species composition ($\chi^2$-based) to a linear combination of environmental
predictors, producing canonical axes that maximise the correlation between
species scores and environmental constraints. The species matrix included taxa
occurring in at least five plots ($n = 200$ species retained); plots with
all-zero species records were excluded ($n = 27$), yielding a 1{,}915~plot
$\times$ 200-species matrix constrained by 24 environmental predictors
(elevation, slope, aspect, \textit{bio1}--\textit{bio19}). The statistical
significance of each canonical axis was assessed by permutation test
(99~permutations). Variance partitioning \citep{borcard1992} was performed to
apportion the adjusted $R^2$ of the full model into unique and shared
fractions attributable to climate variables (\textit{bio1}--\textit{bio19})
and topographic variables (elevation, slope, aspect).

\subsection{Indicator species analysis}
\label{sec:indicator}

Indicator species for each forest type were identified using the IndVal method
\citep{dufresne1997}. IndVal combines specificity ($A$: proportion of a
species' occurrences concentrated in a given group) and fidelity ($B$:
proportion of group plots where the species is present):
\begin{equation}
  \text{IndVal} = A \times B
  \label{eq:indval}
\end{equation}
Significance was assessed by permutation of plot group labels (999
permutations; $\alpha = 0.05$). The minimum attainable $p$-value with 999
permutations is 0.001.

\subsection{Co-occurrence network analysis}
\label{sec:cooccurrence}

Species co-occurrence was examined by constructing a binary
(presence--absence) plot-by-species matrix. Ambiguously labelled records
(taxa matching ``Unknown'' or ``Not listed'') were excluded before network
construction. The co-occurrence matrix was computed as the inner product of
the transposed binary matrix ($\mathbf{B}^\top \cdot \mathbf{B}$), yielding
the number of plots shared by each species pair. An undirected weighted graph
was built with species as nodes; an edge was added between any two species
co-occurring in at least one shared plot (edge weight = shared plot count).
Node-level metrics computed included degree centrality and betweenness
centrality (networkx~$\geq$~3.0). Community structure was detected using the
greedy modularity algorithm \citep{clauset2004}, and modularity $Q$ was
computed as the fraction of edges within communities minus the expected
fraction under a random null model.

To assess whether the observed modularity exceeds chance expectation, a
swap-permutation null model was implemented: each permutation shuffled the
binary matrix using a curveball-type algorithm \citep{strona2014} that
preserves both row (plot richness) and column (species occupancy) marginal
totals. Community detection and $Q$ were computed on each of 99 shuffled
matrices. Null model results are summarised as the standardised effect size
SES $= (Q_{\text{obs}} - \bar{Q}_{\text{null}}) / \text{SD}_{\text{null}}$
and the $p$-value $P(Q_{\text{null}} \geq Q_{\text{obs}})$.

\subsection{Stratification Complexity Index}
\label{sec:sci}

A Stratification Complexity Index (SCI) was computed for each plot to
integrate multi-layer structural information into a single composite score.
Six components were derived: total species richness ($S$), Shannon entropy
($H'$), Simpson diversity ($D$), and layer-specific richness for Trees
($S_\text{T}$), Shrubs ($S_\text{S}$), and Herbs ($S_\text{H}$). Each
component was $z$-standardised across all plots to zero mean and unit
variance. SCI was defined as the unweighted sum of the six $z$-scores:
\begin{equation}
  \text{SCI}_i = z_S + z_{H'} + z_D + z_{S_\text{T}} + z_{S_\text{S}} + z_{S_\text{H}}
  \label{eq:sci}
\end{equation}
The equal-weighting assumption was evaluated by a leave-one-out (LOO)
sensitivity analysis: for each component, a LOO-SCI was computed from the
remaining five, and the Spearman rank correlation between LOO-SCI and full SCI
was recorded. If all correlations exceed $r \geq 0.95$, the equal-weight
assumption is considered empirically supported.

\subsection{Spatially-explicit EVI trend analysis}
\label{sec:evi}

Long-term vegetation greenness trends were assessed using MODIS Terra
Vegetation Indices 16-day composite data (MOD13Q1 Version~061) at native
{\raise.17ex\hbox{$\scriptstyle\sim$}}250~m resolution over the period 2000--2024 (25 annual observations). Three
raster products were exported from Google Earth Engine (GEE):
(i)~Theil-Sen slope raster (pixel-wise robust linear trend estimator
\citep{sen1968}; units = EVI~yr$^{-1}$);
(ii)~Mann-Kendall tau raster with associated $p$-values; where the GEE
$p$-value band contained missing values, analytical $p$-values were derived
from the Mann-Kendall normal approximation
($\text{Var}(S) = n(n-1)(2n+5)/18$;
$Z = (S \mp 1)/\sqrt{\text{Var}(S)}$;
$p = 2(1-\Phi(|Z|))$, with $n = 25$~years); and
(iii)~annual EVI mean stack (25-band composite).

Pixel-level trend classification used the criteria: significant greening
(positive slope, $p \leq 0.05$), significant browning (negative slope,
$p \leq 0.05$), or stable/non-significant. To account for multiple testing
across approximately 700{,}000 simultaneous pixel comparisons,
Benjamini-Hochberg (BH) false discovery rate correction
\citep{benjamini1995} was applied at FDR~$\leq 0.05$. Per-plot EVI slope
and Mann-Kendall tau were extracted at each NFI plot centroid by bilinear
resampling. Stratified summaries of mean Theil-Sen slope and percentage of
plots with significant positive trends were computed by 500-m elevation band
and forest type. The Pearson correlation between per-plot EVI Theil-Sen slope
and species richness was computed; the coefficient of determination ($R^2$) is
reported alongside $r$ to characterise explained variance.

\subsection{Reproducibility and code availability}
\label{sec:reproducibility}

All analyses were implemented in Python~3.12 using a modular, deterministic
pipeline (\texttt{run\_pipeline.py}) with a fixed random seed (seed = 42
throughout). Key dependencies include numpy~$\geq$~1.26, pandas~$\geq$~2.0,
scikit-learn~$\geq$~1.4, scikit-bio~$\geq$~0.6, networkx~$\geq$~3.0,
rasterio~$\geq$~1.3, geopandas~$\geq$~0.14, matplotlib~$\geq$~3.8, and
scipy~$\geq$~1.12. All outputs are fully reproducible from raw inputs given
the fixed seed. Code and documentation are available at [GitHub URL --- to be
inserted at submission].

%% ============================================================
\section{Results}
\label{sec:results}

\subsection{Alpha diversity across elevation and forest type}
\label{sec:res_alpha}

The 1{,}942 NFI plots harboured 2{,}195 unique taxa across the full
elevational range (113--5{,}469~m). Mean species richness across all plots
was 17.1~$\pm$~9.1 taxa per plot (range 1--57). Species richness varied
significantly among the nine 500-m elevation bands
(Kruskal-Wallis $H = 630.8$, $p < 10^{-130}$, $\eta^2 = 0.322$;
\cref{tab:alpha}), confirming that the banding scheme captures a strong and
statistically well-justified gradient.

Richness peaked in the 500--1{,}000~m band
(24.8~$\pm$~9.0 spp per plot; $n = 171$ plots) and declined monotonically to
8.9~$\pm$~5.6 in the highest band ($>$4{,}000~m; $n = 285$ plots;
\cref{tab:alpha}), representing a 2.8-fold reduction across the sampled
gradient. Plots below 500~m ($n = 103$) showed slightly lower richness
(22.9~$\pm$~10.4) than the 500--1{,}000~m peak, consistent with a weak
mid-domain effect at the lowest end of the gradient. Mean Shannon entropy
followed the same trend (overall mean $H' = 1.905~\pm~0.661$; mean
$D = 0.740~\pm~0.190$).

Layer-specific analysis revealed that the tree layer contributed the greatest
share of plot richness (mean 9.1~$\pm$~7.3 spp per plot), followed by the
shrub layer (5.5~$\pm$~4.2) and the herb layer (2.8~$\pm$~3.2). Among the
12 forest types, Subtropical Forest (26.0~$\pm$~8.7 spp per plot) and Warm
Broadleaved Forest (24.2~$\pm$~8.5) supported the highest mean richness, while
Dry Alpine Scrub (8.2~$\pm$~4.7) and Chirpine Forest (11.6~$\pm$~7.1)
supported the lowest.

\subsection{Beta diversity and community turnover}
\label{sec:res_beta}

Two-dimensional NMDS ordination (Bray-Curtis dissimilarity, SMACOF algorithm;
Kruskal stress-1 = 0.338) revealed a strong primary compositional axis (NMDS1,
range $-1.065$ to $+0.968$) separating warm lowland forest types from cold
subalpine communities. Subtropical Forest (NMDS1 mean: $+0.49$) and Warm
Broadleaved Forest ($+0.49$) were strongly separated from
Juniper--Rhododendron Scrub ($-0.62$), Dry Alpine Scrub ($-0.58$), and Fir
Forest ($-0.57$; \cref{fig:nmds}). The Kruskal stress-1 of 0.338 substantially
exceeds the widely accepted threshold of 0.20, indicating that the 2D
projection does not faithfully represent the full Bray-Curtis dissimilarity
structure for this dataset (1{,}942 plots, 2{,}195 species); accordingly, the
ordination is treated as a qualitative visual summary only. A metric PCoA
ordination is provided as Supplementary Fig.~S2.

Compositional differentiation among forest types was confirmed independently
by PERMANOVA, which is not affected by ordination distortion. Forest type
explained 8.9\% of total beta-diversity variation
(pseudo-$F = 17.03$, $R^2 = 0.089$, $p = 0.001$; 999~permutations;
$k = 12$ groups). PERMDISP indicated significant heterogeneity of
within-group dispersions (Kruskal-Wallis $H = 923.09$, $p < 0.001$),
meaning the PERMANOVA signal partly reflects differences in compositional
spread among forest types in addition to centroid separation; both statistics
are reported for transparency.

\subsection{Environmental drivers of community assembly}
\label{sec:res_cca}

Ecological CCA (1{,}915 plots $\times$ 200 species constrained by 24
environmental variables) identified two significant canonical axes (permutation
test, 99 permutations). CCA1 (eigenvalue = 0.873, 2.45\% explained,
$p = 0.01$) captured the dominant temperature--elevation gradient; CCA2
(eigenvalue = 0.435, 1.22\% explained, $p = 0.01$) captured a secondary
precipitation gradient. The full constrained model explained an adjusted
$R^2 = 4.22\%$ of compositional variance (\cref{tab:cca}, \cref{fig:cca}).

The strongest predictor of CCA1 was \textit{bio6} (minimum temperature of the
coldest month, CCA1 score $= +0.963$), followed by \textit{bio11} (mean
temperature of the coldest quarter, $+0.946$), \textit{bio9} (mean
temperature of the driest quarter, $+0.950$), and \textit{bio1} (mean annual
temperature, $+0.940$). Elevation loaded strongly but negatively on CCA1
($-0.945$), confirming that CCA1 represents the thermal gradient from warm
lowland to cold alpine environments. On CCA2, the primary loadings were
\textit{bio18} (precipitation of the warmest quarter, $+0.560$), \textit{bio16}
(precipitation of the wettest quarter, $+0.552$), and \textit{bio12} (mean
annual precipitation, $+0.536$), collectively defining the monsoon-moisture
axis.

Variance partitioning \citep{borcard1992} indicated that the 19 climate
variables accounted for an adjusted $R^2 = 4.01\%$ (pure climate
fraction = 2.90\%), while the three topographic variables accounted for 1.32\%
(pure topography fraction = 0.21\%). The substantial shared fraction
($\approx$1.11\%) reflects the expected collinearity between climate and
elevation in a compressed mountain system.

\subsection{Indicator species by forest zone}
\label{sec:res_indicator}

IndVal analysis identified 416 significant taxon--habitat associations across
the 12 forest types (all $p \leq 0.05$; minimum attainable $p = 0.001$ with
999 permutations; \cref{tab:indval}). The number of significant indicators
varied widely among forest types, from 12 (Fir Forest; Cool Broadleaved
Forest) to 127 (Subtropical Forest). Conifer species exhibited the strongest
habitat fidelity: \textit{Pinus roxburghii} (IndVal = 0.794),
\textit{Picea spinulosa} (0.463), \textit{Pinus wallichiana} (0.411), and
\textit{Abies densa} (0.387). The second-ranked indicator for Chirpine
Forest, \textit{Indigofera dosua} (IndVal = 0.496), is a disturbance-associated
legume shrub ecologically characteristic of the warm, dry, fire-prone
Chirpine belt. Alpine-zone types were characterised by \textit{Rhododendron}
indicators (\textit{R.\ anthopogon}, \textit{R.\ aeruginosum},
\textit{R.\ setosum}), consistent with their known dominance above treeline.

\subsection{Co-occurrence network structure}
\label{sec:res_network}

The species co-occurrence network comprised 1{,}083 nodes (species) and
26{,}875 undirected weighted edges, with a mean degree of 49.6
(\cref{fig:network}). Greedy modularity community detection recovered five
modules (modularity $Q = 0.287$). Against a swap-permutation null model
preserving both species occupancy and plot richness marginals
(99 permutations; $n_{\text{swaps}} = n_{\text{plots}} \times n_{\text{species}}$
per permutation), the observed modularity was not significantly greater than
expected by chance (SES~$= 0.615$, $p = 0.242$), indicating that the
five-module co-occurrence structure does not exceed null-model expectations
at $\alpha = 0.05$. The three largest modules broadly correspond to lowland
broadleaf, montane mixed-forest, and subalpine--alpine assemblages.

Hub species with the highest betweenness centrality among named taxa were
\textit{Maesa chisia} (degree = 550, betweenness = 0.057),
\textit{Ficus}~sp.\ (degree = 510, betweenness = 0.036),
\textit{Rosa sericea} (degree = 249, betweenness = 0.031), and
\textit{Acer campbellii} (degree = 296, betweenness = 0.020).

\subsection{Stratification Complexity Index}
\label{sec:res_sci}

The SCI ranged from $-15.15$ to $+25.15$ across the 1{,}942 plots
(mean = 0 by construction, SD = 5.88). Subtropical Forest
(mean SCI $+4.49~\pm~5.17$) and Warm Broadleaved Forest ($+3.84~\pm~5.13$)
were the most structurally complex forest types; Dry Alpine Scrub
($-5.44~\pm~4.36$) and Juniper--Rhododendron Scrub ($-3.78~\pm~5.06$) were
the least complex (\cref{fig:sci}). Leave-one-out sensitivity analysis
confirmed robust equal-weighting: Spearman correlations between each LOO-SCI
and the full six-component SCI ranged from $r = 0.983$ to $r = 0.999$
(all exceeding the 0.95 threshold), indicating that no single component
dominates the index.

\subsection{Spatially-explicit EVI trend patterns, 2000--2024}
\label{sec:res_evi}

Pixel-level Mann-Kendall trend analysis of the 2000--2024 MODIS EVI record
(25 annual composites; $n = 703{,}900$ valid pixels; 39{,}032~km$^2$ valid
forest area) classified 29.9\% of Bhutan's forest area as significantly
greening after BH FDR correction
(FDR~$\leq 0.05$; 210{,}180~pixels; 11{,}655~km$^2$), while only 0.7\% showed
significant browning (5{,}076~pixels; 282~km$^2$; \cref{tab:evi}, \cref{fig:evimap}).
Using uncorrected nominal $p \leq 0.05$, the greening fraction was 42.4\%
(298{,}245 pixels); both corrected and uncorrected figures are reported in
\cref{tab:evi} for transparency.

Greening intensity was strongly structured by elevation (\cref{fig:evigrad}).
Mean Theil-Sen slope was highest at the lowest elevations
($<$500~m: 0.00187~EVI~yr$^{-1}$; 70.6\% of plots with a significant positive
trend) and declined monotonically through the montane zone to a minimum at
3{,}500--4{,}000~m (0.00061~EVI~yr$^{-1}$; 18.9\% plots significantly
greening), before recovering above 4{,}000~m
(0.00159~EVI~yr$^{-1}$; 43.2\% plots significantly greening), consistent with
upslope shrubification at the alpine treeline ecotone. By forest type,
the highest proportions of positively trending plots were recorded in
Subtropical Forest (94.5\%), Dry Alpine Scrub (94.4\%), and Warm Broadleaved
Forest (94.1\%); the lowest in Spruce Forest (57.1\%) and Blue Pine Forest
(63.0\%).

Per-plot trend classification confirmed: 811 of 1{,}942 plots (41.8\%)
significantly greening ($p \leq 0.05$), 29 plots (1.5\%) significantly
browning, and 1{,}102 plots (56.7\%) stable or non-significant.
EVI Theil-Sen slope correlated positively with per-plot species richness
(Pearson $r = 0.107$, $R^2 = 0.012$, $p = 2.2 \times 10^{-6}$;
\cref{fig:eviscatter}), and more weakly with the SCI
($r = 0.079$, $R^2 = 0.006$, $p = 4.6 \times 10^{-4}$).

%% ============================================================
\section{Discussion}
\label{sec:discussion}

\subsection{Elevation--richness relationship: monotonic decline, not a mid-domain hump}
\label{sec:disc_richness}

Across Bhutan's 1{,}942 NFI plots, species richness peaked at 500--1{,}000~m
(24.8~$\pm$~9.0 spp per plot) and declined monotonically towards the alpine
zone (8.9~$\pm$~5.6 spp above 4{,}000~m), a 2.8-fold reduction driven
primarily by cold-season temperature stress and the truncation of growing-season
length. This pattern contrasts with the hump-shaped mid-domain richness peak
documented in some Himalayan transects \citep{grytnes2002} but is consistent
with studies in wetter, climatically equable mountain systems where the
productivity advantage of lowland environments is not offset by aridity
\citep{rahbek1995}. In Bhutan, the warm, wet subtropical zone benefits from
high precipitation and a long frost-free season that supports complex
multi-layered forests with the highest tree, shrub, and herb layer richness
simultaneously. The slight richness depression at the very lowest elevations
($<$500~m; 22.9~$\pm$~10.4 spp) likely reflects both reduced plot
representation in the narrow foothill zone and the simplifying effect of
land-use pressure near the southern agricultural--forest boundary, where
forest patches are smaller and more disturbed than the contiguous mid-elevation
belt. These dynamics should be tested with repeated inventory cycles that
better capture foothill forest structural heterogeneity.

\subsection{Cold-season temperature as the primary community filter}
\label{sec:disc_cca}

Ecological CCA identified cold-season minimum temperature
(\textit{bio6}, CCA1 score $= +0.963$) and mean temperature of the coldest
quarter (\textit{bio11}, $+0.946$) as the strongest environmental predictors
of floristic composition, jointly defining the primary canonical axis (CCA1:
2.45\% explained, $p = 0.01$). Elevation loaded counter-directionally
($-0.945$), confirming that CCA1 represents the integrated cold-stress
gradient rather than elevation per se. This result is consistent with
\citeauthor{korner2012}'s \citeyearpar{korner2012} synthesis that cold
hardiness and growing-season length are the dominant functional trait filters
structuring tree communities along Himalayan altitudinal gradients. CCA2
(1.22\%, $p = 0.01$) captured a secondary precipitation axis --- monsoon
moisture (\textit{bio18}, $+0.560$) versus rain-shadow dryness ---
distinguishing wet subtropical and warm broadleaved assemblages from dry Blue
Pine and alpine scrub communities. The explained variance of the full model
(adj-$R^2 = 4.22\%$) is low in absolute terms but typical for species-rich
presence--absence matrices \citep{legendre2012}. Variance partitioning
confirmed that the climate fraction substantially outweighs the pure
topographic fraction (pure climate adj-$R^2 = 2.90\%$ vs.\ pure
topography = 0.21\%), and the large shared fraction ($\approx$1.11\%) reflects
the inherent collinearity of temperature and elevation in a compressed mountain
gradient. Critically, these CCA patterns are corroborated by the PERMANOVA
result (pseudo-$F = 17.03$, $R^2 = 0.089$, $p = 0.001$), which is fully
independent of ordination geometry.

\subsection{Conifer fidelity and discrete vertical zonation}
\label{sec:disc_indicator}

The indicator species results confirm that Bhutan's forest types occupy
discrete, identifiable phytogeographic zones rather than forming a single
continuous floristic gradient. Conifers displayed exceptionally high habitat
fidelity --- \textit{Pinus roxburghii} (IndVal = 0.794),
\textit{Picea spinulosa} (0.463), \textit{Pinus wallichiana} (0.411), and
\textit{Abies densa} (0.387) --- reflecting the well-known altitudinal belt
structure of Himalayan conifer forests: Chirpine below {\raise.17ex\hbox{$\scriptstyle\sim$}}1{,}800~m, Blue Pine at
1{,}800--3{,}000~m, Spruce and Hemlock at 2{,}500--3{,}500~m, and Fir above
3{,}000~m \citep{schickhoff2005}. In contrast, broadleaved forest types
produced more numerous but lower-scoring indicators, reflecting their broader
environmental tolerances and higher internal beta-diversity. Alpine zone types
were characterised by \textit{Rhododendron} indicators (\textit{R.\ anthopogon},
\textit{R.\ aeruginosum}, \textit{R.\ setosum}), consistent with the
well-documented dominance of ericaceous scrub above the Himalayan treeline
ecotone. The ecologically meaningful second indicator for Chirpine Forest ---
\textit{Indigofera dosua} (IndVal = 0.496), a disturbance-tolerant legume
shrub --- highlights that indicator analysis captures not only dominant canopy
species but also the characteristic understorey associates that define a forest
type's ecological identity.

\subsection{Network co-occurrence structure and null-model context}
\label{sec:disc_network}

The co-occurrence network (1{,}083 species nodes; 26{,}875 edges; $Q = 0.287$)
recovered five modules, the three largest of which correspond broadly to the
major vegetation belts identified by PERMANOVA and indicator analysis ---
lowland broadleaf, montane mixed-forest, and subalpine--alpine assemblages ---
consistent with elevation-structured co-occurrence patterns. However, the
null-model test (99 swap permutations preserving species occupancy and plot
richness marginals) returned a non-significant result (SES~$= 0.615$,
$p = 0.242$), indicating that the observed modularity does not exceed chance
expectation at $\alpha = 0.05$. While the five modules are spatially
interpretable and ecologically coherent, the co-occurrence structure is not
statistically distinguishable from a random assembly of species sharing the
same occupancy and plot-richness constraints. This finding is not surprising
given that Himalayan species have broad and overlapping elevational ranges
\citep{schickhoff2005}, and the high species richness of the dataset generates
extensive incidental co-occurrence that the null model successfully captures.
The modules should be interpreted as descriptive co-occurrence groupings
rather than evidence of structured assembly processes.

Hub species with high betweenness centrality --- \textit{Maesa chisia}
(betweenness = 0.057), \textit{Ficus}~sp.\ (0.036), \textit{Rosa sericea}
(0.031), \textit{Acer campbellii} (0.020) --- are wide-ranging generalists
present across the lowland-to-montane transition. Their structural centrality
makes them candidates for umbrella or foundation species roles in biodiversity
monitoring: changes in their abundance or range would disproportionately affect
the connectivity of the co-occurrence structure.

\subsection{Stratification Complexity gradient: climate and disturbance}
\label{sec:disc_sci}

The SCI gradient --- from subtropical ($+4.49$) and warm broadleaved ($+3.84$)
at the positive extreme to alpine scrub ($-5.44$) at the negative extreme ---
tracks the combined effect of thermal energy availability and growing-season
length on the development of multi-layered forest structure. High-SCI forests
accrue structural complexity through the simultaneous establishment of dense
tree canopies, diverse shrub understoreys, and rich herb layers --- a process
requiring adequate precipitation and high temperatures throughout an extended
growing season. Two forest types stand out as potential disturbance indicators:
Chirpine Forest (mean SCI~$= -3.01$) is physiognomically simple and
frequently fire-prone, while Non Forest plots (mean SCI~$= -3.13$) represent
degraded or transitional land cover. The LOO sensitivity analysis confirming
$r = 0.983$--$0.999$ for all component substitutions validates SCI as a
robust composite index suitable for cross-inventory standardisation, and
supports its adoption as a structural monitoring metric in future NFI cycles.

\subsection{Spatially-explicit greening: patterns, hypotheses, and constraints}
\label{sec:disc_evi}

The finding that 29.9\% of Bhutan's valid forest area showed significant
positive MODIS EVI trends over 2000--2024 after BH FDR correction
(FDR~$\leq 0.05$; 210{,}180 pixels; 11{,}655~km$^2$; uncorrected nominal
$p \leq 0.05$ figure: 42.4\%) is consistent with the global vegetation
greening phenomenon documented by \citet{zhu2016} using MODIS LAI and NDVI
data, and with Himalayan-specific studies reporting widespread productivity
increases across Nepal, north-east India, and the Tibetan Plateau
\citep{mishra2015, lamsal2017}. The spatial gradient in greening intensity ---
strongest at low elevations ($<$500~m: 70.6\% of plots significantly greening;
mean slope 0.00187~EVI~yr$^{-1}$) and weakest in the montane zone
(3{,}500--4{,}000~m: 18.9\%; slope 0.00061~EVI~yr$^{-1}$), with a secondary
high-elevation recovery ($>$4{,}000~m: 43.2\%; slope
0.00159~EVI~yr$^{-1}$) --- is consistent with at least three non-exclusive
mechanisms. At low elevations, warming-extended growing seasons, CO$_2$
fertilisation, and reduced disturbance under Bhutan's community and religious
forest protection programmes may all contribute to enhanced canopy greenness.
At the highest elevations, the recovery is most likely attributable to upslope
shrubification: the documented advance of \textit{Rhododendron} and
\textit{Juniperus} scrub under atmospheric warming \citep{liang2021}. The
intermediate montane zone showed the weakest greening signal, potentially
reflecting increased cloud cover reducing MODIS signal quality, human
land-use pressure, or slower canopy structural changes in mature forest stands.
These hypotheses cannot be disentangled with a single-epoch inventory and
16-day composite EVI data alone.

The positive but weak correlation between EVI Theil-Sen slope and species
richness ($r = 0.107$, $R^2 = 0.012$, $p = 2.2 \times 10^{-6}$) is
consistent with a productivity-diversity relationship at the landscape scale
but explains only 1.2\% of spatial variation in plot richness, precluding
causal inference. Species diversity in Bhutan's forests is the product of
evolutionary history, dispersal limitation, biotic interactions, and disturbance
legacies that collectively dwarf any productivity signal detectable in 24~years
of EVI data at 250~m resolution.

\subsection{Limitations}
\label{sec:disc_limitations}

Several methodological limitations should be acknowledged. First, the 2D NMDS
Kruskal stress-1 of 0.338 substantially exceeds the conventional acceptability
threshold ($<$0.20), meaning the ordination does not faithfully represent the
Bray-Curtis dissimilarity structure; all beta-diversity inference relies on the
PERMANOVA result, which is geometry-independent; a PCoA is provided as
Supplementary Fig.~S2. Second, the overall explained variance in both PERMANOVA
($R^2 = 0.089$) and CCA (adj-$R^2 = 4.22\%$) is modest, reflecting the
complexity and stochasticity inherent in species-rich tropical-montane
vegetation data and the scale mismatch between plot-level species records and
landscape-scale climate predictors. Third, the NFI represents a single-epoch
survey; temporal trajectories of community change cannot be assessed without
repeat inventories. Fourth, EVI pixel-level significance testing was conducted
using a normal approximation for Mann-Kendall $p$-values, which introduces
minor approximation error at $n = 25$ time points, and the 250~m MODIS pixel
footprint substantially exceeds the scale of individual NFI plots. Fifth,
co-occurrence edges represent spatial co-presence, not demonstrated biotic
interactions. Future work should address these limitations through
repeat-inventory beta-diversity analysis, higher-resolution (Sentinel-2)
vegetation phenology mapping, and mechanistic modelling of community assembly
processes.

\subsection{Management and policy implications}
\label{sec:disc_management}

Four management implications follow directly from these results. First,
subtropical and warm broadleaved forests --- which simultaneously show the
highest species richness, highest SCI, and strongest greening trend --- warrant
the highest biodiversity protection priority in land-use planning and REDD+
accounting frameworks. Second, cold-season temperature thresholds (particularly
the \textit{bio6} and \textit{bio11} gradients identified by CCA) define the
ecologically critical transition zones most vulnerable to climate warming;
monitoring programmes should focus on these thermal boundaries as early-warning
indicators of vegetation zone displacement. Third, the SCI provides a
standardised, repeatable structural complexity metric that can be computed from
any inventory dataset recording species identity by vertical layer, and is
recommended for adoption in future NFI cycles to track long-term forest
structural change. Fourth, the 24-year MODIS EVI greening trend validates the
effectiveness of Bhutan's constitutional forest conservation mandate at the
landscape scale, but monitoring should be disaggregated to forest-type and
elevation levels to detect differential responses and emerging browning signals
--- particularly in Spruce and Blue Pine forests, which showed the lowest
proportions of positively trending plots.

%% ============================================================
\section{Conclusions}
\label{sec:conclusions}

Bhutan's forests display strong vertical zonation structured primarily by
cold-season temperature, with species richness declining 2.8-fold from the
subtropical lowlands to the alpine zone
(Kruskal-Wallis $H = 630.8$, $\eta^2 = 0.322$). PERMANOVA confirmed
significant compositional differentiation among the 12 forest types
(pseudo-$F = 17.03$, $R^2 = 0.089$, $p = 0.001$), and ecological CCA
identified cold-season minimum temperature (\textit{bio6}, score~$= +0.963$)
and elevation ($-0.945$) as the primary community-assembly drivers (CCA1:
2.45\% explained, $p = 0.01$; pure climate adj-$R^2 = 2.90\%$). Indicator
species analysis validated discrete phytogeographic boundaries, with conifers
showing the highest habitat fidelity (\textit{Pinus roxburghii}, IndVal =
0.794). Co-occurrence network analysis (1{,}083 nodes; 26{,}875 edges)
recovered five modules ($Q = 0.287$), the three largest broadly corresponding
to lowland broadleaf, montane mixed, and subalpine--alpine assemblages;
modularity did not significantly exceed null-model expectations
(SES~$= 0.615$, $p = 0.242$), indicating descriptive rather than
statistically structured co-occurrence groupings. The Stratification Complexity
Index captured multi-layer structural differences across all 12 forest types
and was robust to component weighting (all LOO Spearman $r \geq 0.983$),
supporting its adoption as a standardised structural metric for future
inventory cycles. Pixel-level MODIS EVI trend analysis (2000--2024) revealed
that 29.9\% of Bhutan's forested area exhibited significant long-term greening
after BH FDR correction (FDR~$\leq 0.05$; 210{,}180~pixels;
11{,}655~km$^2$), with the strongest signal at low elevations ($<$500~m:
70.6\% of plots significantly greening) and a secondary peak above 4{,}000~m
consistent with alpine shrubification. EVI trend slope correlated positively
with species richness ($r = 0.107$, $R^2 = 0.012$, $p < 0.001$), though the
low explained variance precludes causal inference. Together, these results
provide the first integrated, country-scale quantitative baseline for Bhutan's
forest stratification and establish a fully reproducible open-source framework
for long-term biodiversity and ecosystem monitoring in support of Bhutan's
REDD+ commitments and National Biodiversity Strategy.

%% ============================================================
\section*{CRediT authorship contribution statement}

\textbf{Wangdi}: Conceptualisation, Methodology, Software, Formal analysis,
Data curation, Writing --- original draft, Writing --- review \& editing,
Visualisation, Project administration. \textbf{Sangay Chedup}: Data curation,
Field data collection, Writing --- review \& editing.

\section*{Declaration of competing interest}

The authors declare that they have no known competing financial interests or
personal relationships that could have appeared to influence the work reported
in this paper.

\section*{Data availability}

The National Forest Inventory data are held by the Department of Forests and
Park Services, Ministry of Energy and Natural Resources, Royal Government of
Bhutan. Requests for access should be directed to the corresponding author.
The analysis pipeline, configuration, and all derived outputs are available at
[GitHub URL --- to be inserted at submission] under an open-source licence.

\section*{Acknowledgements}

The authors thank the NFI field teams of the Department of Forests and Park
Services for data collection, and acknowledge the support of the Ministry of
Energy and Natural Resources, Royal Government of Bhutan.
[Additional funding sources and institutional acknowledgements to be completed
by lead author.]

%% ============================================================
%% TABLES
%% ============================================================

\clearpage

\begin{table}[htbp]
\caption{Alpha diversity summary by 500-m elevation band.
Kruskal-Wallis $H = 630.8$, $p < 10^{-130}$, $\eta^2 = 0.322$.
Full forest-type breakdown in Supplementary Table~S3.}
\label{tab:alpha}
\centering
\begin{tabular}{lrrrr}
\toprule
Elevation band & $n$ plots & Mean richness & SD & Mean $H'$ \\
\midrule
$<$500 m        & 103 & 22.9 & 10.4 & -- \\
500--1{,}000 m  & 171 & \textbf{24.8} & 9.0  & -- \\
1{,}000--1{,}500 m & 249 & 21.0 & 9.5  & -- \\
1{,}500--2{,}000 m & 255 & 21.6 & 9.3  & -- \\
2{,}000--2{,}500 m & 258 & 17.6 & 7.2  & -- \\
2{,}500--3{,}000 m & 253 & 16.0 & 6.0  & -- \\
3{,}000--3{,}500 m & 204 & 14.2 & 5.6  & -- \\
3{,}500--4{,}000 m & 164 & 11.1 & 5.8  & -- \\
$>$4{,}000 m    & 285 &  8.9 & 5.6  & -- \\
\midrule
\textbf{All plots} & \textbf{1{,}942} & \textbf{17.1} & \textbf{9.1} & \textbf{1.905} \\
\bottomrule
\end{tabular}
\end{table}

\clearpage

\begin{table}[htbp]
\caption{Top three indicator species per forest zone (IndVal $= A \times B$;
sorted by IndVal; all $p \leq 0.05$, 999 permutations, minimum attainable
$p = 0.001$). Full list of all 416 significant indicators in Supplementary
Table~S2.}
\label{tab:indval}
\centering
\small
\begin{tabular}{lrp{3.2cm}rp{3.2cm}r}
\toprule
Forest zone & $n$ sig. & Top indicator & IndVal & 2nd indicator & IndVal \\
\midrule
Chirpine Forest         & 26  & \textit{Pinus roxburghii}        & \textbf{0.794} & \textit{Indigofera dosua}          & 0.496 \\
Spruce Forest           & 44  & \textit{Picea spinulosa}          & 0.463          & \textit{Pieris formosa}            & 0.264 \\
Blue Pine Forest        & 36  & \textit{Pinus wallichiana}        & 0.411          & \textit{Elsholtzia fruticosa}      & 0.245 \\
Fir Forest              & 12  & \textit{Abies densa}              & 0.387          & \textit{Rhododendron hodgsonii}    & 0.107 \\
Subtropical Forest      & 127 & \textit{Pterospermum acerifolium} & 0.314          & \textit{Tabernaemontana divaricata}& 0.293 \\
Evergreen Oak Forest    & 22  & \textit{Quercus lamellosa}        & 0.301          & \textit{Persea}~sp.                & 0.208 \\
Hemlock Forest          & 20  & \textit{Tsuga dumosa}             & 0.286          & \textit{Betula utilis}             & 0.165 \\
Dry Alpine Scrub        & 46  & \textit{Rhododendron anthopogon}  & 0.283          & \textit{R.\ aeruginosum}           & 0.166 \\
Warm Broadleaved Forest & 45  & \textit{Maesa chisia}             & 0.173          & \textit{Syzygium}~sp.              & 0.135 \\
Juniper Rhod.\ Scrub    & 13  & \textit{Juniperus recurva}        & 0.102          & \textit{Rhododendron lepidotum}    & 0.086 \\
Cool Broadleaved Forest & 12  & \textit{Elatostema}~sp.           & 0.089          & \textit{Aconogonon molle}          & 0.051 \\
Non Forest              & 13  & \textit{Saussurea}~sp.            & 0.081          & \textit{Bistorta macrophylla}      & 0.053 \\
\bottomrule
\end{tabular}
\end{table}

\clearpage

\begin{table}[htbp]
\caption{CCA environmental biplot scores for top-ranked variables on CCA1
(temperature--elevation gradient) and CCA2 (precipitation gradient).
Full-model adjusted $R^2 = 4.22\%$; both axes significant ($p = 0.01$,
99~permutations).}
\label{tab:cca}
\centering
\begin{tabular}{llrrl}
\toprule
Variable & Description & CCA1 score & CCA2 score & Primary axis \\
\midrule
\textit{bio6}  & Min.\ temp.\ coldest month     & $+0.963$ & --       & Temperature \\
\textit{bio9}  & Mean temp.\ driest quarter     & $+0.950$ & --       & Temperature \\
\textit{bio11} & Mean temp.\ coldest quarter    & $+0.946$ & --       & Temperature \\
Elevation      & m a.s.l.                       & $-0.945$ & --       & Temperature \\
\textit{bio1}  & Mean annual temperature        & $+0.940$ & --       & Temperature \\
\textit{bio5}  & Max.\ temp.\ warmest month     & $+0.926$ & --       & Temperature \\
\textit{bio4}  & Temperature seasonality        & $-0.866$ & --       & Temperature \\
\textit{bio7}  & Temperature annual range       & $-0.794$ & --       & Temperature \\
\textit{bio18} & Precip.\ warmest quarter       & --       & $+0.560$ & Precipitation \\
\textit{bio16} & Precip.\ wettest quarter       & --       & $+0.552$ & Precipitation \\
\textit{bio13} & Precip.\ wettest month         & --       & $+0.548$ & Precipitation \\
\textit{bio12} & Annual precipitation           & --       & $+0.536$ & Precipitation \\
\textit{bio3}  & Isothermality                  & --       & $+0.525$ & Precipitation \\
\bottomrule
\end{tabular}
\end{table}

\clearpage

\begin{table}[htbp]
\caption{Pixel-level MODIS EVI trend classification (MOD13Q1, 2000--2024;
Mann-Kendall). BH = Benjamini-Hochberg FDR correction applied at
FDR~$\leq 0.05$ to account for {\raise.17ex\hbox{$\scriptstyle\sim$}}700{,}000 simultaneous pixel comparisons.
Both corrected and uncorrected figures are reported for transparency.}
\label{tab:evi}
\centering
\begin{tabular}{lrrrr}
\toprule
Trend class & $n$ pixels & Area (km$^2$) & \% valid area \\
\midrule
\multicolumn{4}{l}{\textit{BH FDR-corrected (FDR $\leq$ 0.05) --- preferred}} \\
\quad Significant greening & 210{,}180 & 11{,}655 & \textbf{29.9\%} \\
\quad Stable / non-significant & 488{,}644 & 27{,}096 & 69.4\% \\
\quad Significant browning  & 5{,}076   & 282      & 0.7\%  \\
\midrule
\multicolumn{4}{l}{\textit{Uncorrected (nominal $p \leq$ 0.05)}} \\
\quad Significant greening & 298{,}245 & 16{,}538 & 42.4\% \\
\quad Stable / non-significant & 396{,}621 & 21{,}993 & 56.4\% \\
\quad Significant browning  & 9{,}034   & 501      & 1.3\%  \\
\midrule
\textbf{Total valid} & \textbf{703{,}900} & \textbf{39{,}032} & 100\% \\
\bottomrule
\end{tabular}
\end{table}

%% ============================================================
%% FIGURE CAPTIONS
%% ============================================================

\clearpage
\section*{Figure captions}

\begin{figure}[htbp]
  \centering
  % \includegraphics[width=\textwidth]{figs/map_species_richness.png}
  \caption{\textbf{Study area and NFI plot distribution.}
    Bhutan (26.7$^\circ$--28.4$^\circ$N, 88.7$^\circ$--92.2$^\circ$E)
    with 1{,}942 NFI plots coloured by forest type. Elevation is shown
    by background shading. Inset: location of Bhutan within the Eastern
    Himalayas.}
  \label{fig:map}
\end{figure}

\begin{figure}[htbp]
  \centering
  % \includegraphics[width=\textwidth]{figs/alpha_diversity_elevation.png}
  \caption{\textbf{Alpha diversity across the elevational gradient.}
    Per-plot species richness ($S$), Shannon entropy ($H'$), and
    layer-specific richness (Trees, Shrubs, Herbs) as a function of
    elevation. Solid lines: LOESS smoother; shaded bands: 95\% CI.
    Overall Kruskal-Wallis $H = 630.8$, $p < 10^{-130}$, $\eta^2 = 0.322$.}
  \label{fig:alpha}
\end{figure}

\begin{figure}[htbp]
  \centering
  % \includegraphics[width=\textwidth]{figs/nmds_ordination.png}
  \caption{\textbf{NMDS ordination of plot species composition.}
    Two-dimensional NMDS (Bray-Curtis dissimilarity; SMACOF algorithm;
    Kruskal stress-1 = 0.338). Points coloured by forest type. The high
    stress value indicates that the 2D projection does not faithfully
    represent Bray-Curtis distances; this figure is a qualitative visual
    summary only. Inferential statistics (PERMANOVA:
    pseudo-$F = 17.03$, $R^2 = 0.089$, $p = 0.001$) are independent of
    ordination geometry. A metric PCoA is provided as Supplementary
    Fig.~S2.}
  \label{fig:nmds}
\end{figure}

\begin{figure}[htbp]
  \centering
  % \includegraphics[width=\textwidth]{figs/cca_sites.png}
  \caption{\textbf{Ecological CCA triplot of species, sites, and
    environmental vectors.}
    Sites (grey points), environmental arrows (red, scaled by CCA1/CCA2
    score), and selected species scores (blue labels) for the two
    significant canonical axes (CCA1: 2.45\%, CCA2: 1.22\%;
    both $p = 0.01$, 99 permutations; full model adj-$R^2 = 4.22\%$).
    Arrow length proportional to CCA score.}
  \label{fig:cca}
\end{figure}

\begin{figure}[htbp]
  \centering
  % \includegraphics[width=\textwidth]{figs/species_co_occurrence_network.png}
  \caption{\textbf{Species co-occurrence network.}
    1{,}083 species nodes (coloured by module) and 26{,}875 edges
    (edge weight = shared plot count). Node size proportional to degree.
    Hub species (highest betweenness centrality) labelled.
    Modularity $Q = 0.287$; null-model SES $= 0.615$, $p = 0.242$
    (not significantly different from null).}
  \label{fig:network}
\end{figure}

\begin{figure}[htbp]
  \centering
  % \includegraphics[width=\textwidth]{figs/map_sci_index.png}
  \caption{\textbf{Stratification Complexity Index (SCI).}
    (a) Spatial distribution of SCI across all 1{,}942 NFI plots
    (RdYlGn colour scale; centred at 0). (b) Mean SCI by forest type
    (error bars: $\pm$1~SD). SCI $= z_S + z_{H'} + z_D +
    z_{S_\text{T}} + z_{S_\text{S}} + z_{S_\text{H}}$;
    all leave-one-out Spearman $r \geq 0.983$.}
  \label{fig:sci}
\end{figure}

\begin{figure}[htbp]
  \centering
  % \includegraphics[width=\textwidth]{figs/evi_spatial_trend_map.png}
  \caption{\textbf{Spatially-explicit MODIS EVI trend map, 2000--2024.}
    Pixel-level Mann-Kendall trend classification at {\raise.17ex\hbox{$\scriptstyle\sim$}}250~m resolution
    after BH FDR correction (FDR~$\leq 0.05$): green = significant
    greening (29.9\%; 210{,}180 pixels; 11{,}655~km$^2$), brown =
    significant browning (0.7\%), grey = stable. NFI plot centroids
    shown as white circles.}
  \label{fig:evimap}
\end{figure}

\begin{figure}[htbp]
  \centering
  % \includegraphics[width=0.9\textwidth]{figs/evi_slope_vs_elevation.png}
  \caption{\textbf{EVI Theil-Sen slope by elevation zone and forest type.}
    (a) Per-plot Theil-Sen slope (EVI~yr$^{-1}$) vs.\ elevation with
    rolling mean $\pm$ SE ribbon. Elevation zone bands shaded.
    (b) Box plots of per-plot EVI slope by forest type, sorted by median;
    $n$ per group annotated.}
  \label{fig:evigrad}
\end{figure}

\begin{figure}[htbp]
  \centering
  % \includegraphics[width=0.9\textwidth]{figs/evi_slope_vs_richness_sci.png}
  \caption{\textbf{EVI Theil-Sen slope vs.\ species richness and SCI.}
    Per-plot EVI slope ($y$-axis) vs.\ (a) species richness and
    (b) Stratification Complexity Index ($x$-axis), coloured by elevation
    (terrain\_r palette). OLS regression line with bootstrap 95\% CI band.
    Richness: $r = 0.107$, $R^2 = 0.012$, $p < 0.001$;
    SCI: $r = 0.079$, $R^2 = 0.006$, $p < 0.001$.}
  \label{fig:eviscatter}
\end{figure}

%% ============================================================
%% REFERENCES
%% ============================================================

\bibliographystyle{elsarticle-num-names}
\bibliography{references}

\end{document}
